\section{Introduction}\label{sec:introduction}
An economic decision often requires a continuation value that must itself be computed. A household changes consumption and portfolio risk while investing in its preferences. A sophisticated consumer anticipates future selves, and a firm invests against the policies of its rivals. These problems differ in their utility domains and deviation criteria, but each requires an evaluation of the consequences that follow a proposed action. The preference-formation problem of \citet{qi2025} provides a motivating example in which internal and external choices interact.

This paper develops Neural Bellman Operators (NBO), an evaluation and improvement map for these economic problems. A critic estimates the continuation value of a specified policy. A feasible actor uses its derivative channels to change current decisions, keeping the evaluated continuation fixed. The resulting policy is assessed against an explicit economic criterion. Boundary conditions, information, numerical approximation, and the deployed decision rule are part of that criterion.

The central question is when learning a continuation value makes economic computation more effective. A fitted continuation can be reused across states and action queries, but fitting is costly and its approximation error can alter the selected action. A credible comparison must allow a simulation-based method to cache the information that the same decision problem permits. It must also include the cost of training and choosing the continuation representation. We therefore study the accuracy and complete work of the learned evaluation block when several economic decisions use a common continuation. The nonlinear capital economy provides the common primitives for this comparison and for the policy-specific error analysis.

The theory follows the evaluation--decision--payoff chain. A fixed-policy equation identifies the continuation being approximated. The improvement map specifies the actor's feasible action and the derivatives it can use. General verification results distinguish an evaluation defect from unfinished action maximization. For the capital economy, an exact finite Bellman identity places the relevant continuation error on the action change under the candidate's own occupation law. A signed continuation calculation can assess the same identity using common future paths, retaining the cancellation between predicted improvement and prediction error. Its payoff range depends on the economy and action tube. The continuous-time account separately charges the implementation, diffusion, clipping, and numerical errors of the held controller. These are error bounds for computed objects; they do not assert a convergence rate for a neural optimizer.

The computational design uses two complementary forms of comparison. A repeated-decision study examines the incremental value of reusing a learned continuation, including feasible Raw caching and the complete cost of fitting and selection. Independent confirmation on the original capital policies examines the previously reported paired payoff comparisons using a proof, constants, reporting code, and confidence allocation fixed before the new simulation bank. A strengthened neural-HJB implementation uses an exact diffusion trace, independent residual selection, and a documented action implementation. A separate signed Bellman experiment assesses the precision of the occupation account. The design and reporting preserve the distinction between new decision problems, new payoff observations on existing policies, and changes in the analytical treatment of existing observations.

The continuation menu gives a direct economic result. In the two previously unpiloted fifty-dimensional Long and Intermediate designs, NBO's protected payoff gains over both the materialized Raw actor and direct policy exceed the prespecified margin $10^{-4}$. Its complete construction and query time is also lower than direct policy's in those cells. At the final prescribed budget, six of eight cells have NBO--cached-SAA intervals wholly within $\pm10^{-4}$, with lower construction and query work for NBO. Cached SAA delivers materially higher payoffs in the other two cells, and the cheaper earlier SAA budgets remain in the comparison. These results concern the complete sixteen-stream populations and 384 fixed task-state queries in each finite economy. They identify useful fitted-continuation decisions and their observed costs. A separate assessment below tests continuation prediction directly, and a simultaneous catalogue calculation compares every declared budget; neither establishes a unique causal mechanism or a minimum-work frontier over all possible procedures.

The independent original-policy confirmation also makes the paired comparison prospective: NBO exceeds the original neural-HJB implementation by $10^{-4}$ in both dimensions. A newly trained, strengthened HJB procedure gives a more demanding comparison across four continuous-diffusion economies. Only the low-volatility, fifty-dimensional cell establishes that margin against both its greedy and distilled deployments; the other direct comparisons remain unresolved. The signed Bellman assessment gives continuous-economy gain lower bounds above $0.0005$ in both dimensions while its signed continuation-correction intervals still contain zero. The full adapted-class regret bounds remain distinct from these gains and from the finite menu's scalar accuracy certificates.

The previous complete sixteen-stream experiment remains part of the current evidence. Its direct NBO--Raw and NBO--DPO payoff intervals were unresolved, and cached Raw was cheaper at the fifty-dimensional online target. A common-innovation transfer refinement made the NBO--HJB comparison positive in both ten and fifty dimensions, but the refinement was developed after the original analytical freeze and inspection of interval widths. The original dimension-ten interval included zero. We retain both intervals and their chronology. The original mechanism lower bound was nonpositive in each dimension. Those observations motivate the current experiments and remain visible alongside them; improving a verifier does not by itself identify a benefit of learning the continuation.

Two further results connect this evidence to the learned object and the economic accuracy target. Independent prediction assessment compares the nonlinear scalar continuation with realized Raw caches at the same selected action pairs. Common future-path noise cancels from the difference of squared errors. Eight of 24 prespecified intervals establish lower NBO contrast risk, four favor Raw, and twelve remain unresolved. This measures a conditional prediction benefit without inferring payoff ordering from risk ordering. Separately, exact closure of all 216 simultaneous menu payoff intervals bounds regret relative to the complete sixteen-candidate budget catalogue. In the fifty-dimensional Intermediate design, final-stage NBO is the least-cost candidate certified within $10^{-4}$ of the best catalogue payoff. Its mean complete construction and query cost is 12.6993 seconds, rounded upward. Cheaper uncertified candidates remain possible competitors, and the full cost of conducting the comparison remains charged. This last calculation is an explicitly post-freeze deterministic use of simultaneous intervals, not a new prospective stopping experiment.

The economic applications remain part of the paper's substantive scope. Recursive utility evaluates each policy through its own utility process. Endogenous preference adjustment connects an internal shadow value with consumption and portfolio demand. Temporal selves require a local spike-deviation equilibrium, whereas strategic firms require a complete unilateral best response against fixed rivals. The main article states these economic distinctions and their principal implications. The current \emph{Economic Applications} companion contains their full models, derivations, proof arguments, and numerical comparisons. The technical supplement contains the general theory, detailed capital accounts, and complete computational records. This organization places the common numerical question in the foreground while preserving the distinct economic problems that give the operator its content.
